% D5 CJ uniform proof progress 057

\begin{proposition}[Carry\_jump reduction to the flat corner]
Let
\[
B=(s,u,v,\lambda,f),\qquad \tau,\qquad \epsilon_4\in\{\mathrm{flat},\mathrm{wrap},\mathrm{carry\_jump},\mathrm{other}\},
\]
and let \(\rho=u_{\mathrm{source}}+1\) be used only as an auxiliary proof device.
Define
\[
\delta:=\rho-(s+u+v+\lambda)\in\mathbb Z_m.
\]
Assume the checked identities
\[
q\equiv u-\rho+1 \pmod m
\quad\text{on }\epsilon_4=\mathrm{carry\_jump},
\qquad
q\equiv 1-s-v-\lambda \pmod m
\quad\text{on }\epsilon_4=\mathrm{carry\_jump}.
\]
Then on a carry\_jump boundary state \(x\) with \(c(x)=0\), its successor
\(y=F(x)\) is flat and satisfies
\[
\delta(y)=m-2,
\qquad
s(y)=s(x)+1.
\]
Consequently, if the flat-corner rule
\[
\epsilon_4(Fy)=
\begin{cases}
\mathrm{other},& \delta(y)=m-2 \text{ and } s(y)=2,\\
\mathrm{flat},& \delta(y)=m-2 \text{ and } s(y)\neq 2,
\end{cases}
\]
holds uniformly in odd \(m\), then the carry\_jump reset law follows:
\[
R_{\mathrm{cj}}(x)=
\begin{cases}
0,& q(x)=m-1,\\
1,& q(x)\neq m-1 \text{ and } s(x)=1,\\
m-2,& q(x)\neq m-1 \text{ and } s(x)\neq 1.
\end{cases}
\]
\end{proposition}

\begin{proof}[Proof sketch]
The zero-reset branch is exactly the carry slice \(q=m-1\).  On the noncarry
branch, carry\_jump has \(dn=(1,1,0,0)\), so \(s\) and \(u\) both increase by one
while \(v\) and \(\lambda\) remain fixed.  The formula for \(q\) on the
carry\_jump boundary implies that \(q\) is unchanged across the step, and hence
\[
\delta(y)=\rho-(s(y)+u(y)+v(y)+\lambda(y))
=\rho-(s(x)+u(x)+v(x)+\lambda(x)+2)
\equiv -2\equiv m-2.
\]
Thus every noncarry carry\_jump state lands in a flat state with
\(\delta=m-2\).  If the flat-corner rule holds, then this successor has
\(\tau=1\) exactly when \(s(y)=2\), and otherwise \(\tau=m-2\).  Since
\(s(y)=s(x)+1\), this is equivalent to the stated noncarry dichotomy.
\end{proof}
